\chapter{Experimental Methodology}
\label{ch:experiments}
\label{ch:methodology}

\section{Research Design}
The study uses a fixed-case comparative experimental design. All controllers
run on the same rig, firmware measurement pipeline, actuator interface, safety
logic, and logger. The independent factor is the selected controller mode.
The observed responses are pendulum angle, arm angle, estimated angular
rates, commanded acceleration, commanded velocity, trial outcome, and
diagnostic events.

This is not a repeated-measures statistical experiment. The manifest contains
one representative trial for each available controller--condition pair, and
the conclusions are intentionally limited to those trials.

\section{Dataset definition}
All results in this thesis are generated from one fixed dataset of experimental runs selected through the manifest \texttt{experiments/manifest\_thesis\_20260209.json}. This makes the analysis repeatable: the same thesis assets can be rebuilt from the same seven pinned runs.

\subsection*{Thesis run identifiers}
For readability, the internal run labels used by the analysis pipeline are mapped to concise thesis identifiers:
\begin{itemize}
  \item Hold-at-center: \textbf{H1--H3}
  \item Commanded reference tracking (nudge): \textbf{N1}
  \item Finger tap disturbance: \textbf{T1--T3}
\end{itemize}
Each condition is represented by one pinned run only, so later comparisons are qualitative and metric-based rather than statistical. Table~\ref{tab:dataset_manifest} lists the run identifiers used in this thesis.

\InputRequired{tables/tab_dataset_manifest.tex}

\section{Controller and pipeline configuration}
The controller and pipeline settings were fixed for the pinned runs and were
captured in each session event log. The complete parameter inventory is
reported in Appendix~\ref{app:implementation_parameters}; the controller
chapter reports the smaller mode-specific subsets needed to interpret each
control law.

\section{Calibration and Trial Procedure}
Before each run, the operator disarmed the system, recorded the active
configuration, selected the controller, and calibrated the arm at its center
with the pendulum held upright. The system was then armed and allowed to
engage automatically only when the pendulum was inside the upright and
low-rate engagement window. At the end of a run, the system was explicitly
disarmed.

The operational sequence was:
\[
\texttt{!}\ \rightarrow\ \texttt{G}\ \rightarrow\ \texttt{C 0/1/2}\
\rightarrow\ \texttt{Z}\ \rightarrow\ \texttt{E}\ \rightarrow\ \texttt{!}.
\]
Sign settings had previously been checked through the diagnostic procedure
described in Section~\ref{ch:notation}. The pinned runs used
\texttt{motorSign=1}, \texttt{ALPHA\_SIGN=-1},
\texttt{THETA\_SIGN=1}, and \texttt{CTRL\_SIGN=1}.

\section{Experimental Conditions}
\begin{itemize}
  \item \textbf{Hold}: engage near $\theta_{\mathrm{ref}}=0$ and observe
  undisturbed upright behavior.
  \item \textbf{Nudge}: command a sequence of bounded arm-angle targets with
  the linear controller while maintaining upright balance.
  \item \textbf{Tap}: apply brief manual disturbances while the controller is
  balancing. Tap onset and magnitude are not independently measured.
\end{itemize}

\section{Logging artifacts}
Each logger session produces:
\begin{itemize}
  \item a device-time, nine-field time-series log (about \SI{50}{Hz}),
  \item a host-time event log (status and diagnostic messages).
\end{itemize}
Within a session, each controller engagement followed by a disarm defines one \emph{trial}.

The hold and nudge runs are straightforward to interpret from the recorded trajectories. The tap runs are less direct: they do not contain explicit disturbance markers, so the dataset supports qualitative disturbance-recovery comparisons but not precise impulse-response timing metrics.

\section{Event-to-CSV alignment and limitations}
The event log timestamps are host wall-clock time, while the time-series log uses device time. Therefore, naive alignment can introduce \SIrange{100}{500}{ms} shifts and corrupt any time-window metrics.

The analysis pipeline aligns host events to device time by matching sparse status snapshots to nearby device-time samples using a multi-signal error metric (for example $\alpha$, $\dot{\alpha}$, $\theta$, $\dot{\theta}$, acceleration, and velocity) under a monotonicity constraint. A linear host-to-device mapping is then fit and residuals are reported. If alignment quality is unreliable, the pipeline avoids precise event-time annotations and instead reports metrics computed directly from device-time trajectories.

\section{Metric definitions}
All metrics are defined once here and reused in later tables:
\begin{itemize}
  \item \textbf{RMS and max:} $\mathrm{RMS}(|\alpha|)$, $\max|\alpha|$, $\max|\theta|$, and related quantities over the trial duration.
  \item \textbf{Drift slope:} slope of $\theta(t)$ over the trial (least-squares fit), reported in \si{\degree\per\second}.
  \item \textbf{Effort:} RMS and max of the commanded base acceleration $u$ in \si{\step\per\second\squared}.
  \item \textbf{Saturation and clamp percentages:} fraction of samples where acceleration saturation is active and fraction where the base is near the mechanical limit.
\end{itemize}

For the nudge trial, rise time, overshoot, steady-state error, and entry into a
fixed $\pm\SI{1}{\degree}$ settling band are evaluated for each commanded
step. A missing settling time means the criterion was not maintained within
the available post-command window.

\section{Reproducibility and Traceability}
The analysis command resolves the manifest into explicit session paths and
writes trial-level metrics, summary tables, alignment diagnostics, figures,
and a resolved manifest. The thesis export copies the required raw metric
tables into \texttt{thesis/tables/raw/} and records the analysis report name,
manifest name, repository revision, and export time in
\texttt{analysis\_metadata.tex}. This allows reported values to be traced to
the selected sessions without treating ignored raw logs as thesis source
files.
